\chapter{INTRODUCTION}
\section{Background of the Study}
Construction rework denotes an activity undertaken to correct, replace, remove or repeat work already executed. Its consequences differ from those of an ordinary construction task: the previous work and the corrective intervention must both be established before resources or delay can be attributed to rework. Research on construction incidents and project management has shown that errors, changes, coordination and workmanship can interact, but reported effects vary by case and context \parencite{forcada2017,love2022,asadi2023}. A defect or crack is therefore a possible indicator for investigation rather than a verified rework event.

The construction of reinforced-concrete (RCC) buildings includes numerous associated activities whose failure can require replacement of finishes, repairs to adjacent work or further intervention after occupation. Three technically connected areas warrant attention in the present study: waterproofing at bathrooms and terraces; plain-cement-concrete (PCC) pavement associated with building sites; and plaster cracking linked in part to masonry workmanship. These are construction-process subjects, not mutually exclusive descriptions of completed rework events. For example, a waterproofing failure can require removing previously installed tiles, while a masonry junction defect can become visible only after plastering.

This study uses a practitioner account and a separate photographic evidence collection. In subsequent clarification the researcher identified himself as the contributor to the practitioner account and reported direct involvement in several Pokhara construction projects. These statements are respondent-reported, and the contemporaneous site diaries and project records still require verification. The original mixed Nepali--English account and an edited technical translation are preserved in the project research-input record dated October 2026. Attribution to the researcher is now self-reported, while the original site identities, dates and event-to-document mapping are not independently established in the current manuscript. Accordingly, this thesis treats the supplied material as \emph{reported field observations}, not a sampled dataset or independently verified local case series. The reported conditions offer an applied basis for examining plausible mechanisms and needed site records; established technical references provide a separate basis for engineering interpretation.

\section{Statement of the Problem}
A site team may recognize damp patches next to bathrooms, cracks in concrete access pavements or plaster cracks near intersecting walls. These visible conditions do not uniquely establish their causes. Water may enter through an incomplete wet-area membrane, a penetration or another source; a pavement crack may develop because of restraint, curing, joint layout or support; a wall-junction crack may reflect differential movement, inadequate bonding or another mechanism. Without distinguishing what was observed from what has been diagnosed, a report may incorrectly blame a worker, recommend the wrong corrective action or count routine maintenance as rework.

Previous rework studies emphasize incident records, causal relationships, cost and prevention \parencite{liu2020,garg2021,mostofi2022}. Their empirical estimates do not establish what happened in an individual Pokhara project. The present research problem is to interpret three groups of locally reported quality concerns using sound construction engineering, and to identify which supplementary observations and documents would be needed to establish actual corrective rework and its consequences.

\section{Research Questions}
\begin{enumerate}
\item What waterproofing, PCC pavement and masonry/plaster quality concerns have been reported by a practitioner in relation to construction work in Pokhara?
\item Which construction mechanisms and relevant technical requirements plausibly explain these concerns, and which explanations remain unverified?
\item How should observed conditions, proven corrective work, contributing causes and resource or schedule consequences be distinguished and documented?
\item What preventive inspection and workmanship measures are supported by the technical literature and applicable project requirements?
\end{enumerate}

\section{Objectives of the Study}
\subsection{General Objective}
To assess practitioner-reported quality concerns and potential rework mechanisms related to waterproofing, PCC pavement and plaster/masonry work associated with RCC building construction in Pokhara Metropolitan City, using a transparent separation of observations, technical explanation and verified corrective action.

\subsection{Specific Objectives}
\begin{enumerate}
\item To document and classify the three categories of field concerns without creating unsupported project-level counts or causal findings.
\item To interpret reported conditions against the ten established construction-rework papers and independently identifiable engineering guidance.
\item To set out the evidence needed to confirm corrective work and attribute additional materials, labour, costs or activity delays.
\item To formulate realistic construction-quality and rework-prevention recommendations for each category.
\end{enumerate}

\section{Study Area and Scope}
The geographical focus is Pokhara Metropolitan City. The research title retains its focus on RCC building construction projects. Waterproofing and plaster works are considered as associated building activities. PCC pavements are considered only where they form part of an eligible building project's access road, internal circulation area, courtyard or directly related site-development works. A separate municipal or highway pavement is not automatically a case within this thesis; such work would require scope authorization and suitable records.

\begin{table}[htbp]
\caption{Three main construction subjects and admissible study boundary}
\label{tab:worktypes}
\centering
\begin{tabularx}{\textwidth}{@{}L{1.0cm}L{3.6cm}X@{}}
\toprule Code & Subject & Relationship to the building-project research scope\\ \midrule
W01 & Waterproofing & Bathrooms, wet rooms, terraces and relevant penetrations, junctions or finishes.\\
W02 & PCC pavement & Plain concrete access and circulation pavement demonstrably belonging to an eligible RCC building project.\\
W03 & Plaster cracking / masonry & Plaster defects considered with brick bonds, junctions, mortar preparation, curing and required wall detailing.\\ \bottomrule
\end{tabularx}
\end{table}

\section{Evidence and Limitations}
The preserved account and its translation are professional narratives. Twenty-five photographs have subsequently been inventoried, but their photo-to-site assignments, corresponding original diaries, drawings, crack measurements, repair instructions and cost records have not yet been independently reconciled as study evidence. An approximate professional experience duration does not constitute a sample size. The present work therefore reports qualitative practitioner concerns and an engineering interpretation rather than frequencies, statistically ranked causes or monetary losses.

An observed condition is not sufficient proof of rework. Verification requires both previous execution and subsequent corrective action. The technical origin of the condition and its consequences are distinct issues, each of which may remain undetermined even when a repair is documented. Applicable code requirements are evaluated according to building type, approved design and the actual specification; examples from international manuals are not presumed to be compulsory Nepalese rules.

\section{Organization of the Thesis}
Chapter 2 reviews construction-rework research and technical guidance for the three subjects. Chapter 3 explains the qualitative interpretation and case-verification framework. Chapter 4 analyses the supplied practical concerns and corresponding corrective and preventive measures. Chapter 5 presents conclusions and recommendations within the evidence available. The appendices provide a structured case-recording form and photograph/evidence categories for future corroboration.
